% Part: first-order-logic
% Chapter: proof-systems
% Section: axiomatic-deduction

\documentclass[../../../include/open-logic-section]{subfiles}

\begin{document}

\iftag{FOL}
      {\olfileid[pt]{fol}{prf}{axd}}
      {\olfileid[pt]{pl}{prf}{axd}}

\olsection{Dedução Axiomática}

As !!{derivation}s axiomáticas são os sistemas de !!{derivation}
lógicos mais antigos e mais simples. Suas !!{derivation}s são
simplesmente sequências de !!{sentence}s.  Uma sequência de
!!{sentence}s conta como uma !!{derivation} correta se toda
!!{sentence}~$!A$ nela satisfaz uma das seguintes condições:
\begin{enumerate}
\item $!A$ é um axioma, ou
\item $!A$ é !!a{element} de um dado conjunto~$\Gamma$ de !!{sentence}s, ou
\item $!A$ é justificada por uma regra de inferência.
\end{enumerate}
Para ser um axioma, $!A$ tem que ter a forma de um entre vários
esquemas fixos de !!{sentence}. Há muitos conjuntos de esquemas de
axiomas que fornecem um sistema de !!{derivation} satisfatório
(correto e completo) para a lógica de primeira ordem. Alguns são
organizados de acordo com os conectivos que governam, por exemplo, os
esquemas
\[
!A \lif (!B \lif !A) \qquad !B \lif (!B \lor !C) \qquad (!B \land !C) \lif !B
\]
são axiomas comuns que governam $\lif$, $\lor$ e~$\land$. Alguns
sistemas de axiomas visam um número mínimo de axiomas. Dependendo dos
conectivos tomados como primitivos, é até possível encontrar sistemas
de axiomas que consistem em um único axioma.

Uma regra de inferência é uma afirmação condicional que dá uma
condição suficiente para que uma !!{sentence} em uma !!{derivation}
seja justificada. O modus ponens é uma regra muito comum desse tipo:
diz que se $!A$ e $!A \lif !B$ já estão justificados, então $!B$ está
justificado. Isso significa que uma linha em uma !!{derivation} que
contém a !!{sentence}~$!B$ está justificada, desde que tanto $!A$
quanto $!A \lif !B$ (para alguma !!{sentence}~$!A$) apareçam na
!!{derivation} antes de~$!B$.

A relação $\Proves$ baseada em !!{derivation}s axiomáticas é definida
da seguinte forma: $\Gamma \Proves !A$ se, e somente se, há uma
!!{derivation} com a !!{sentence}~$!A$ como sua última fórmula (e
$\Gamma$ é tomado como o conjunto de !!{sentence}s nessa !!{derivation}
que são justificadas por~(2) acima).  $!A$ é um teorema se~$!A$ tem
uma !!{derivation} em que~$\Gamma$ é vazio, isto é, toda !!{sentence}
na !!{derivation} é justificada por (1) ou~(3). Por exemplo, aqui está
uma !!{derivation} que mostra que $\Proves !A  \lif (!B \lif (!B \lor
!A))$:
\begin{derivation}
  1. & $!B \lif (!B \lor !A)$ \\
  2. & $(!B \lif (!B \lor !A)) \lif (!A  \lif (!B \lif (!B \lor !A)))$\\
  3. & $!A  \lif (!B \lif (!B \lor !A))$
\end{derivation}
A !!{sentence} na linha~1 tem a forma do axioma $!A \lif (!A \lor
!B)$ (com os papéis de $!A$ e $!B$ invertidos). A sentença na
linha~2 tem a forma do axioma $!A \lif (!B \lif !A)$. Assim, ambas
as linhas estão justificadas. A linha~3 é justificada por modus
ponens: se a abreviarmos como $!D$, então a linha~2 tem a forma $!C
\lif !D$, onde $!C$ é $!B \lif (!B \lor !A)$, isto é, a linha~1.

Um conjunto $\Gamma$ é inconsistente se $\Gamma \Proves \lfalse$. Um
sistema de axiomas completo também provará que $\lfalse \lif !A$ para
qualquer~$!A$, e assim, se $\Gamma$ é inconsistente, então $\Gamma
\Proves !A$ para qualquer~$!A$.

Sistemas de !!{derivation}s axiomáticas para lógica foram dados pela
primeira vez por Gottlob Frege em sua \emph{Begriffsschrift} de 1879,
que por essa razão é frequentemente considerada a primeira obra da
lógica moderna. Foram aperfeiçoados em \emph{Principia Mathematica} de
Alfred North Whitehead e Bertrand Russell e por David Hilbert e seus
alunos na década de 1920. São assim frequentemente chamados de
``sistemas de Frege'' ou ``sistemas de Hilbert''. São muito versáteis
no sentido de que é frequentemente fácil encontrar um sistema axiomático
para uma lógica. Porque as !!{derivation}s têm uma estrutura muito
simples e apenas uma ou duas regras de inferência, também é
relativamente fácil provar coisas \emph{sobre} eles. No entanto, são
muito difíceis de usar na prática, isto é, é difícil encontrar e escrever
provas.

\end{document}
